\begin{textsample}{POS Dim 3 – pb – Score 20.00 – pb/pb\_inf\_11.txt}  \label{ex:f3_pos_088}
3.5.1.5.
CORPO, \textbf{GESTOS} E MOVIMENTO Toda experiência da \textbf{criança} acontece por meio do corpo, que ganha \textbf{centralidade} no processo de aprendizagem e desenvolvimento.
É por meio do seu corpo que a \textbf{criança} conhece o mundo, experimenta sensações, expressa sentimentos e emoções, faz escolhas, toma decisões e aproxima -se de sua cultura.
Ao \textbf{manipular} objetos de diferentes formas, cores, \textbf{volumes}, \textbf{pesos} e \textbf{texturas}, ela passa a conhecer e identificar suas sensações, suas funções, potencialidades e limites.
A exploração do ambiente e do espaço por meio do movimento da \textbf{criança} pode ser vivenciada nas áreas internas e externas da \textbf{instituição}, considerando a potência de experiências que envolvem o \textbf{lúdico}.
Elas devem ser ricas e desafiadoras e possibilitar a ampliação e o aprimoramento do movimento e do conhecimento acerca de si e do outro.
Vivências relacionadas a jogos, circuitos, danças, \textbf{brincadeiras} com o movimento e deslocamento do corpo, dentre outras, são atraentes para a \textbf{criança}.
Compete à \textbf{instituição} de Educação \textbf{Infantil} promover situações significativas de aprendizagem, explorando a música, o teatro, as \textbf{brincadeiras}, a dança, o faz de conta para que a \textbf{criança} produza conhecimento sobre o universo cultural e social ampliando o repertório a partir de vivências com movimentos.
Além disso, nessas experiências e em muitas outras, as \textbf{crianças} também se \textbf{deparam}, frequentemente, com conhecimentos matemáticos ( contagem, ordenação, relações entre quantidades, dimensões, medidas, comparação de \textbf{pesos} e de comprimentos, avaliação de distâncias, reconhecimento de formas geométricas, conhecimento e reconhecimento de numerais cardinais e ordinais etc. ) que igualmente aguçam a \textbf{curiosidade}.
Portanto, a Educação \textbf{Infantil} precisa promover experiências nas quais as \textbf{crianças} possam fazer observações, \textbf{manipular} objetos, investigar e explorar seu entorno, levantar hipóteses e consultar fontes de informação para buscar respostas às suas \textbf{curiosidades} e indagações.
Assim, a \textbf{instituição} escolar está criando oportunidades para que as \textbf{crianças} ampliem seus conhecimentos do mundo físico e sociocultural e possam utilizá -los em seu cotidiano.
Com o corpo ( por meio dos sentidos, \textbf{gestos}, movimentos impulsivos ou intencionais, coordenados ou espontâneos ), as \textbf{crianças}, desde cedo, exploram o mundo, o espaço e os objetos do seu entorno, estabelecem relações, expressam -se, \textbf{brincam} e produzem conhecimentos sobre si, sobre o outro, sobre o universo social e cultural, tornando -se, progressivamente, conscientes dessa corporeidade.
Por meio das diferentes linguagens, como a música, a dança, o teatro, as \textbf{brincadeiras} de faz de conta, elas se comunicam e se expressam no entrelaçamento entre corpo, emoção e linguagem.
As \textbf{crianças} conhecem e reconhecem as sensações e funções de seu corpo e, com seus \textbf{gestos} e movimentos, identificam suas potencialidades e seus limites, desenvolvendo, ao mesmo tempo, a consciência sobre o que é seguro e o que pode ser um risco à sua integridade física.
Na Educação \textbf{Infantil}, o corpo das \textbf{crianças} ganha \textbf{centralidade}, pois ele é o \textbf{partícipe} privilegiado das práticas pedagógicas de \textbf{cuidado} físico, orientadas para a emancipação e a liberdade, e não para a submissão.
Assim, a \textbf{instituição} escolar precisa promover oportunidades ricas para que as \textbf{crianças} possam sempre ser animadas pelo espírito \textbf{lúdico} e na interação com seus \textbf{pares}, explorar e vivenciar um amplo repertório de movimentos, \textbf{gestos}, olhares, \textbf{sons} e \textbf{mímicas} com o corpo, para descobrir variados modos de ocupação e uso do espaço com o corpo ( tais como sentar com apoio, rastejar, engatinhar, escorregar, caminhar \textbf{apoiando} -se em berços, mesas e cordas, saltar, escalar, equilibrar -se, correr, dar cambalhotas, alongar -se etc. ).

% matched lemmas: apoiar, brincadeira, brincar, centralidade, criança, cuidado, curiosidade, deparar, gesto, infantil, instituição, lúdico, manipular, mímico, par, partícipe, peso, som, textura, volume
\end{textsample}
